% Introduction scaffold (Phase 2): five named-paragraph stubs covering the
% BHL four-stage scaffold (genre_conventions_theory.md §(c)). Stage 4 is
% split into two paragraphs — the contribution announcement and a
% one-sentence-per-section structural map — following the PSALTer §I
% "Aims" + "Structure" convention. No subsections — letter-style §1.
% Target ~900 words across the five paragraphs.
%
% Paragraph titles are content-bearing micro-headings matching the BHL
% letter idiom (compare 2406.12826: "Introduction to torsion", "No-go
% theorem for vector torsion", "In this letter"). Avoid stage labels.

\section{Introduction}%
\label{Introduction}

% Stage 1 — Frontier tension (~200 w).
% Two-source framing per the recent style refinement:
%   (a) GW detector precision era probes the gauge-gravity sector;
%   (b) the Gertsenshtein literature itself is contested (Hwang–Noh,
%       Palessandro tensions; resolved in §2.3 EM-conventions paragraph).
% Required citations: ≥1 future detector (\cite{amaro2017lisa},
% \cite{nanograv2023fifteen}) and ≥1 current bound (\cite{abbott2023gwtc3}).
\paragraph*{High-frequency gravitational waves and the gauge-gravity sector}
% TODO

% Stage 2 — Deep problem (~250 w).
% Narrows to the deep problem with two moves:
%  - The PGT theory space is multi-dimensional; manual Gertsenshtein-channel
%    derivation is infeasible at scale across this space.
%  - Ghost-freedom imposes a sharp restriction (only spin-1 trace torsion
%    can propagate ghost-free in PGT — Lin–Hobson–Lasenby 2019; Barker–Marzo
%    type analysis), so the Gertsenshtein channel structure inside the
%    ghost-free subspace is a quantitative open question, while the
%    behaviour outside that subspace is itself informative.
% Required citations: \cite{lin2019pgt}, \cite{blagojevic2018hamiltonian};
% optional: 2507.05349 if added to bib. The ghost-restriction logic is the
% conceptual hinge of §1; do not omit. The bridging move that connects
% this to the §4 sector survey (which spans both ghost-restricted and
% non-ghost-restricted theories) is critical: nulls outside the
% ghost-free subspace are evidence FOR the restriction, not a separate
% inquiry.
\paragraph*{Torsion theory space and the spin-1 ghost restriction}
% TODO

% Stage 3 — Prior derivations and their tensions (~250 w).
% Dense citation cluster covering:
%  - Original Gertsenshtein/Boccaletti/Raffelt
%    (\cite{gertsenshtein1962wave}, \cite{boccaletti1970conversion},
%     \cite{raffelt1988mixing});
%  - BHL PGT landscape (2406.12826, 2101.02645, conformal embedding);
%  - PSALTer/Hamilcar as adjacent automation tools
%    (\cite{barker2024psalter}, \cite{barker2025hamilcar});
%  - Modern Gertsenshtein derivations and their tensions
%    (\cite{hwangnoh2023graviton}, \cite{hwangnoh2024nonlinear},
%     \cite{palessandro2024gertsenshtein}, Domcke 2301.02072,
%     Dandoy 2406.17853).
% Close with: "However, no systematic Gertsenshtein channel decomposition
% for PGT+EM has been performed."
\paragraph*{Prior derivations and their tensions}
% TODO

% Stage 4a — Contribution announcement (~150 w).
% Must be specific enough that an examiner reading only §1 knows what to
% expect from §4. The slot enumerates the actual outcomes of the survey
% rather than promising abstract capability. State, in order:
%  (i)   what TIDAL is (one sentence — see par:TidalScope in §3 for the
%        PSALTer/Hamilcar contrast);
%  (ii)  the survey perimeter (which PGT theory classes were scanned and
%        the parameter dimensions of each — match §4 "Survey perimeter"
%        paragraph);
%  (iii) the validation result (TIDAL reproduces the Boccaletti formula);
%  (iv)  the structural finding visible in the kinetic matrix
%        (\cref{HxAxIdentity});
%  (v)   the headline survey finding from the campaign as it stands at
%        submission time (left flexible — campaigns ongoing at scaffold
%        time).
\paragraph*{In this report}
% TODO

% Stage 4b — Structural map (~50 w).
% One sentence per remaining section. Per-section drafting fills in the
% specific verb choice ("In §2 we...", "§3 describes...", etc.).
\paragraph*{Structure of this report}
% TODO
